% Section 5.1, from lectures/L05/appendix/a_function_templates.md.

\section{Function Templates}
\label{c5:sec:function}\label{c5:app:a}
Templates allow functions to operate on multiple data types without duplicating code.

In embedded systems, templates are often used to implement generic utilities without run-time
overhead, while still allowing the compiler to optimize the code for each specific type.

\subsection{Basic example}
\label{c5:sec:basic}
Consider a simple function template that adds two values:

\begin{cppcode}
/**
 * @brief Add two numbers.
 *
 * @tparam T The numeric type.
 *
 * @param[in] x The first value.
 * @param[in] y The second value.
 *
 * @return The sum of the two values.
 */
template<typename T>
constexpr T add(const T x, const T y) noexcept
{
    return x + y;
}
\end{cppcode}

\noindent The function can be used with different types:

\begin{cppcode}
const auto sum1 = add(1, 2);
const auto sum2 = add(1.0, 2.0);
\end{cppcode}

\noindent The compiler generates a separate version of the function for each type.

\subsection{Embedded example: bit manipulation}
\label{c5:sec:bitexample}
In embedded systems, templates are often used for low-level utilities.

The example below sets a bit in a register of arbitrary integral type:

\begin{cppcode}
#include <cstdint>

/**
 * @brief Set bit in register.
 *
 * @tparam T The register type.
 *
 * @param[in, out] reg Register to write to.
 * @param[in] bit The bit to set.
 */
template<typename T>
constexpr void set(T& reg, const std::uint8_t bit) noexcept
{
    reg |= (static_cast<T>(1U) << bit);
}
\end{cppcode}

\noindent Example usage:

\begin{cppcode}
std::uint8_t reg{};
set(reg, 1U);
set(reg, 2U);
\end{cppcode}

\noindent This allows the same function to be reused for different register sizes.

\subsection{Template instantiation}
\label{c5:sec:instantiation}
Templates are instantiated at compile time.

For example:

\begin{cppcode}
std::uint8_t reg1{};
std::uint32_t reg2{};

set(reg1, 1U);
set(reg2, 12U);
\end{cppcode}

\noindent The compiler generates two separate functions:

\begin{cppcode}
set(std::uint8_t&, std::uint8_t)
set(std::uint32_t&, std::uint8_t)
\end{cppcode}

\noindent Each version is compiled independently.

\subsection{Code size considerations}
\label{c5:sec:codesize}
Templates are often described as zero-cost abstractions:
\begin{itemize}
  \item No run-time overhead.
  \item Code is specialized for each type.
\end{itemize}
However, in embedded systems this also means:
\begin{itemize}
  \item Multiple instantiations $\rightarrow$ increased binary size.
  \item More types $\rightarrow$ more generated code.
\end{itemize}
Therefore:
\begin{itemize}
  \item Use templates when they provide clear value.
  \item Avoid unnecessary instantiations in resource-constrained systems.
\end{itemize}

\subsection{Function overloading vs templates}
\label{c5:sec:overloading}
C++ supports function overloading:

\begin{cppcode}
constexpr void set(std::uint8_t& reg, std::uint8_t bit) noexcept;
constexpr void set(std::uint32_t& reg, std::uint8_t bit) noexcept;
\end{cppcode}

\noindent Templates provide a more scalable solution:

\begin{cppcode}
template<typename T>
constexpr void set(T& reg, std::uint8_t bit) noexcept;
\end{cppcode}

\subsection{Summary}
\label{c5:sec:functionsummary}
\begin{itemize}
  \item Function templates enable generic and reusable code.
  \item Templates are instantiated at compile time.
  \item Each type results in a separate generated function.
  \item Templates provide no run-time overhead, but may increase code size.
  \item Templates are especially useful for:
  \begin{itemize}
    \item Bit manipulation.
    \item Utility functions.
    \item Hardware abstraction helpers.
  \end{itemize}
\end{itemize}
